

%% 
%% This is a LaTeX document generated by automatic conversion of a
%% notebook-format document using Publicon 1.0.
%% 

\documentclass{amsart}

\usepackage{amscd,amssymb,amsmath,graphicx,verbatim}
\usepackage[dvips]{hyperref}
\usepackage{textcomp}
\usepackage[OT1,T1]{fontenc}

\newtheorem{theorem}{Theorem}[section]
\newtheorem{lemma}[theorem]{Lemma}
\newtheorem{corollary}[theorem]{Corollary}
\newtheorem{proposition}[theorem]{Proposition}
\newtheorem{conjecture}[theorem]{Conjecture}

\theoremstyle{definition}
\newtheorem{definition}[theorem]{Definition}
\newtheorem{example}[theorem]{Example}
\newtheorem{axiom}[theorem]{Axiom}
\newtheorem{question}[theorem]{Question}

\theoremstyle{remark}
\newtheorem{remark}[theorem]{Remark}

\numberwithin{equation}{section}


\begin{document}


\title{ The core of adjoint functors}
\author{Ross Street}
\address{Centre of Australian Category Theory, Department of Mathematics,
Macquarie University, NSW 2109, Australia}
\curraddr{Department of Mathematics, Macquarie University, NSW 2109,
Australia}
\email{ross.street@mq.edu.au}
\urladdr{http://www.maths.mq.edu.au/{\textasciitilde}street/}\label{I1}
\date{1 December 2011}
\subjclass{18A40; 18D20}
\keywords{adjoint functor; enriched category; counit.}
\begin{abstract}
There is a lot of redundancy in the usual definition of adjoint
functors. We define and prove the core of what is required.\ \ \ 
\end{abstract}
\maketitle

Kan \cite{KanAdjFun} introduced the notion of adjoint functors.
By defining the unit and counit natural transformations, he paved
the way for the notion to be internalized to any 2-category. This
was done by Kelly \cite{Kelly1969} whose interest at the time was
particularly in the 2-category of $\mathcal{V}$-categories for a
monoidal category $\mathcal{V}$ (in the sense of Eilenberg-Kelly
\cite{EilKel1966}). 


For $\mathcal{V}$-categories $\mathcal{A}$ and $\mathcal{X}$, an
{\itshape adjunction} consists of
\begin{enumerate}
\item  $\mathcal{V}$-functors $U:\mathcal{A}\longrightarrow \mathcal{X}$
and $F:\mathcal{X}\longrightarrow \mathcal{A}$;
\item a $\mathcal{V}$-natural family of isomorphisms $\pi :\mathcal{A}(
F X,A) \cong \mathcal{X}( X,U A) $ in $\mathcal{V}$ indexed by $A\in
\mathcal{A}$, $X\in \mathcal{X}$.
\end{enumerate}
We write $\pi :F\dashv U:\mathcal{A}\longrightarrow \mathcal{X}$.


The following result is well known; for example see Section 1.11
of \cite{KellyBook}.
\begin{proposition}

Suppose $U:\mathcal{A}\longrightarrow \mathcal{X}$ is a $\mathcal{V}$-functor,
$F:\mathit{ob}\mathcal{X}\longrightarrow \mathit{ob}\mathcal{A}$
is a function, and, for each $X\in \mathcal{X}$, $\pi :\mathcal{A}(
F X,A) \cong \mathcal{X}( X,U A) $ is a family of isomorphisms $\mathcal{V}$-natural
in $A\in \mathcal{A}$. Then there exists a unique adjunction $\pi
:F\dashv U:\mathcal{A}\longrightarrow \mathcal{X}$ for which $F:\mathrm{ob\mathcal{X}}\longrightarrow
\mathrm{ob\mathcal{A}}$ is the effect of the $\mathcal{V}$-functor
$F:\mathcal{X}\longrightarrow \mathcal{A}$ on objects.\ \ \ \ \ \ \label{XRef-Proposition-121173012}
\end{proposition}

\begin{definition}

For $\mathcal{V}$-categories $\mathcal{A}$ and $\mathcal{X}$, an
{\itshape adjunction core} consists of
\begin{enumerate}
\item functions $U:\mathit{ob}\mathcal{A}\longrightarrow \mathit{ob}\mathcal{X}$
and $F:\mathit{ob}\mathcal{X}\longrightarrow \mathit{ob}\mathcal{A}$;\ \ 
\item a family of isomorphisms $\pi :\mathcal{A}( F X,A) \cong \mathcal{X}(
X,U A) $ in $\mathcal{V}$ indexed by $A\in \mathcal{A}$, $X\in \mathcal{X}$.
\end{enumerate}
\end{definition}


Given such a core, we make the following definitions:
\begin{enumerate}
\item  $\eta _{X}:X\longrightarrow U F X$ is the composite
\[
I\overset{j}{\longrightarrow }\mathcal{A}( F X,F X) \overset{\pi
}{\longrightarrow }\mathcal{X}( X,U F X) ;
\]
\item  $\varepsilon _{A}:F U A\longrightarrow A$ is the composite
\[
I\overset{j}{\longrightarrow }\mathcal{X}( U A,U A) \overset{\pi
^{-1}}{\longrightarrow }\mathcal{A}( F U A,X) ;
\]
\item  $U_{A B}:\mathcal{A}( A,B) \longrightarrow \mathcal{X}( U
A,U B) $ is the composite
\[
\mathcal{A}( A,B) \overset{\mathcal{A}( \varepsilon _{A},1) \ \ }{\longrightarrow
}\mathcal{A}( F U A,B) \overset{\pi }{\longrightarrow }\mathcal{X}(
U A,U B) ;
\]
\item  $F_{X Y}:\mathcal{X}( X,Y) \longrightarrow \mathcal{A}( F
X,F Y) $ is the composite
\[
\mathcal{X}( X,Y) \overset{\mathcal{X}( 1,\eta _{Y}) \ \ }{\longrightarrow
}\mathcal{X}( X,U F Y) \overset{\pi ^{-1}}{\longrightarrow }\mathcal{A}(
F X,F Y) .
\]
\end{enumerate}



Clearly each adjunction $\pi :F\dashv U:\mathcal{A}\longrightarrow
\mathcal{X}$ includes an adjunction core as part of its data. Then
it follows directly from the Yoneda lemma and the definitions (a)
and (b) that the effect of $U$ and $F$ on homs are as in (c) and
(d).

\begin{theorem}

An adjunction core extends to an adjunction if and only if one of
the diagrams (\ref{XRef-Equation-121162956} or (\ref{XRef-Equation-121163025})
below commutes. The adjunction is unique when it exists.
\begin{gather}
\begin{array}{ccc}
 \mathcal{A}( A,B) \otimes \mathcal{A}( F X,A)  & \overset{U_{A
B}\otimes \pi \ \ }{\longrightarrow } & \mathcal{X}( U A,U B) \otimes
\mathcal{X}( X,U A)  \\
 \left. {}{}^{\mathit{comp}}\right\downarrow   & = & \downarrow
^{\mathit{comp}} \\
 \mathcal{A}( F X,B)  & \operatorname*{\longrightarrow }\limits_{\ \ \pi
\ \ } & \mathcal{X}( X,U B) 
\end{array}%
\label{XRef-Equation-121162956}
\end{gather}
\begin{equation}
\begin{array}{ccc}
 \mathcal{X}( Y,U A) \otimes \mathcal{X}( X,Y)  & \overset{\pi ^{-1}\otimes
F_{X Y}\ \ }{\longrightarrow } & \mathcal{A}( F Y,A) \otimes \mathcal{A}(
F X,F Y)  \\
 \left. {}{}^{\mathit{comp}}\right\downarrow   & = & \downarrow
^{\mathit{comp}} \\
 \mathcal{X}( X,U A)  & \operatorname*{\longrightarrow }\limits_{\ \ \pi
^{-1}} & \mathcal{A}( F X,A) 
\end{array}%
\label{XRef-Equation-121163025}
\end{equation}
\end{theorem}

{\itshape Proof} We deal first with the version involving diagram
(\ref{XRef-Equation-121162956}). For an adjunction, (\ref{XRef-Equation-121162956})
expresses the $\mathcal{V}$-naturality of $\pi $ in $A\in \mathcal{A}$.
Conversely, given an adjunction core satisfying (\ref{XRef-Equation-121162956}),
we paste to the left of (\ref{XRef-Equation-121162956} with $X =U
C$, the diagram 
\begin{equation}
\begin{array}{ccc}
 \mathcal{A}( A,B) \otimes \mathcal{A}( C,A)  & \overset{1\otimes
\mathcal{A}( \varepsilon _{C},1) \ \ }{\longrightarrow } & \mathcal{A}(
A,B) \otimes \mathcal{A}( F U C,A)  \\
 \left. {}{}^{\mathit{comp}}\right\downarrow   & = & \downarrow
^{\mathit{comp}} \\
 \mathcal{A}( C,B)  & \operatorname*{\longrightarrow }\limits_{\ \ \mathcal{A}(
\varepsilon _{C},1) \ \ } & \mathcal{A}( F U C,B) 
\end{array}
\end{equation}
which commutes by naturality of composition. This leads to the following
commutative square.
\begin{equation}
\begin{array}{ccc}
 \mathcal{A}( A,B) \otimes \mathcal{A}( C,A)  & \overset{U\otimes
U\ \ }{\longrightarrow } & \mathcal{X}( U A,U B) \otimes \mathcal{X}(
U C,U A)  \\
 \left. {}{}^{\mathit{comp}}\right\downarrow   & = & \downarrow
^{\mathit{comp}} \\
 \mathcal{A}( C,B)  & \operatorname*{\longrightarrow }\limits_{\ \ U\ \ }
& \mathcal{X}( U C,U B) 
\end{array}%
\label{XRef-Equation-121172723}
\end{equation}
We also have the equality
\begin{equation}
\left( I\overset{j}{\longrightarrow }\mathcal{A}( A,A) \overset{U}{\longrightarrow
}\mathcal{X}( U A,U B) \right) =\left( I\overset{j}{\longrightarrow
}\mathcal{X}( U A,U B) \right) %
\label{XRef-Equation-121172748}
\end{equation}
straight from the definitions (b) and (c). Together (\ref{XRef-Equation-121172723})
and (\ref{XRef-Equation-121172748}) tell us that $U$ is a $\mathcal{V}$-functor.
Now the general diagram (\ref{XRef-Equation-121162956}) expresses
the $\mathcal{V}$-naturality of $\pi $ in $A$. By Proposition \ref{XRef-Proposition-121173012},
we have an adjunction determined uniquely by the core.

Writing $\mathcal{V}^{\mathrm{rev}}$ for $\mathcal{V}$ with the
reversed monoidal structure $A\otimes ^{\mathrm{rev}}B=B\otimes
A$, and applying the first part of this proof to the $\mathcal{V}^{\mathrm{rev}}$-enriched
adjunction $\pi ^{-1}:U^{\mathrm{op}}\dashv F^{\mathrm{op}}:\mathcal{X}^{\mathrm{op}}\longrightarrow
\mathcal{A}^{\mathrm{op}}$, which is the same as an adjunction $\pi
:F\dashv U:\mathcal{A}\longrightarrow \mathcal{X}$, we see that
it is equivalent to an adjunction core satisfying (\ref{XRef-Equation-121163025}).\ \ 

\begin{corollary}

If $\mathcal{V}$ is a poset then adjunction cores are adjunctions.
\end{corollary}

{\itshape Proof} All diagrams, including (\ref{XRef-Equation-121162956}),
commute in such a $\mathcal{V}$.\ \ 

\appendix
\begin{thebibliography}{000}
\bibitem{KanAdjFun} D. M. Kan, Adjoint functors, Transactions of
the American Mathematical Society \textbf{87} (1958), 294--329.\label{KanAdjFun}
\bibitem{Kelly1969} G. M. Kelly, Adjunction for enriched categories,
Lecture Notes in Mathematics \textbf{106} (Springer-Verlag, 1969),
166--177.\label{Kelly1969}
\bibitem{EilKel1966} S. Eilenberg and G. M. Kelly, Closed categories,
\textit{Proceedings of the Conference on Categorical Algebra (La
Jolla, 1965)}, Springer-Verlag, Berlin-New York, 1966, 421--526.\label{EilKel1966}
\bibitem{KellyBook} G. M. Kelly, \textit{Basic Concepts of Enriched
Category Theory}, Cambridge University Press, Lecture Notes in Mathematics
64, Cambridge, 1982. \href{http://www.tac.mta.ca/tac/reprints/articles/10/tr10abs.html}{http://www.tac.mta.ca/tac/reprints/articles/10/tr10abs.html}
\label{KellyBook}
\end{thebibliography}

\end{document}
